In decentralized finance (DeFi), lenders can offer flash loans to
borrowers, i.e., loans that are only valid within a blockchain transaction
and must be repaid with some fees by the end of that transaction. Unlike
normal loans, flash loans allow borrowers to borrow large assets
without upfront collaterals deposits. Malicious adversaries use flash loans
to gather large assets to launch costly 
exploitation targeting DeFi protocols. 

In this thesis, we introduce a new framework for automated synthesis of
adversarial transactions that exploit DeFi protocols using flash loans. To
bypass the complexity of a DeFi protocol, we propose a new technique to
approximate the DeFi protocol functional behaviors using numerical methods 
(polynomial linear regression and nearest-neighbor interpolation). 
We then construct an optimization query using the approximated functions 
of the DeFi protocol to find an adversarial attack constituted of a sequence of 
functions invocations with optimal parameters that gives the maximum profit. 
To improve the accuracy of the approximation, we propose a new 
counterexample-driven approximation refinement technique. 
We implement our framework in a tool named {\name}. 
We evaluate {\name} on $16$ DeFi protocols that were
victims to flash loan attacks and $2$ DeFi protocols from Damn Vulnerable DeFi
challenges. {\name} automatically synthesizes an adversarial attack for
$16$ of them. 